\begin{claim}{D.55}{Proved}{Full primary faithfulness and the nine finite torus types}
Every positive involution $b$ in the full
once-punctured-torus loop-outer group has
\[
                      A_p(b)=-I
                           \quad\text{for every prime }p.
                                                               \tag{D.55a}
\]
Consequently, for each prime $p$, the actual
homology map $A_p$ is injective on every
finite subgroup of that group.
The possible finite abstract subgroup
types are exactly the following nine,
all realized:
\[
       1,\ C_2,\ C_3,\ C_4,\ C_2\times C_2,\ C_6,\
                                      S_3,\ D_8,\ D_{12}.
                                                               \tag{D.55b}
\]
The groups $D_8,D_{12}$ have orders eight
and twelve respectively.

For an arbitrary order-four element $a$,
\[
                A_p(a)^2=-I\quad\text{for every }p.
                                                               \tag{D.55c}
\]
At a negative multiplier coordinate $q$
it is necessarily a scalar matrix
$i_q I$ or $-i_q I$, where $i_q^2=-1$
in $\Q_q$; thus $q\equiv1\pmod4$.
Together with (D.54b), this retains
the precise scalar cubic and quartic
multiplier cases. The abstract
classification (D.55b) does not
settle those scalar cases or the
conjugacy of arbitrary finite subgroups.
\end{claim}
\begin{proof}
\emph{The positive involution's full cusp normalizer.}
Let $b$ be any positive involution.
D.43 gives its global multiplier $1$.
D.27 gives an actual representative
$\beta$ of order two, and put
$E_b=F\rtimes_\beta\langle s\rangle$.
The first half of D.46 applies to
every such positive involution; it
does not require $b$ to have first
been exhibited as an order-four square.
It gives
\[
        A_2(b)=-I,\qquad
        H^1(C_F(S),\F_2)=0
          \quad\text{for every order-two section }S.
                                                               \tag{D.55d}
\]

We recall the prime-independent part
of the closed-subgroup argument in
D.48. If $C\leq F$ is closed and
$H^1(C,\F_p)=0$, it cannot contain
the full subgroup $c^{e_p\Z_p}$ of
the torus cusp, where $c=[x,y]$.
If the $p$-part of $[F:C]$ were
finite, choose an open neighborhood
$W\geq C$ with maximal such part.
All subsequent indices are prime
to $p$, so transfer makes the
restriction maps on first cohomology
injective. Their direct limit is
$H^1(C,\F_p)$, contradicting the
nonzero first cohomology of the open
full free group $W$.
If the $p$-part is infinite, choose
an open $W\geq C$ of index divisible
by $p$. Its selected cusp lift has
degree prime to $p$, because it
contains the full cusp $\Z_p$.
There must be at least two boundary
lifts: a unique lift would have
degree $[F:W]$.
A boundary character of this finite
cover then restricts nontrivially
to that cusp $\Z_p$ and hence to
$C$, again a contradiction.
These are the actual open-subgroup
restriction and finite-cover maps
used in D.48, with an arbitrary
prime $p$ in place of three.

Apply this with $p=2$ to (D.55d).
An involution in $N_{E_b}(\langle c\rangle)$
would be a section and would centralize
the entire cusp, since the global
multiplier is $1$.
Its full fixed group would therefore
contain the cusp $\Z_2$, which has
just been excluded.
Thus this cusp normalizer is
torsion-free.
Choose its labelled element $v$:
\[
              vcv^{-1}=c,\qquad v^2=c^\kappa.
                                                               \tag{D.55e}
\]
If $\kappa_2$ were even, the equation
$2z=-\kappa$ would be soluble in
$\widehat\Z$: at two by that
divisibility, and at every odd prime
because two is a unit.
Then $(vc^z)^2=1$, and its nontrivial
quotient label would make it an
involution in the cusp normalizer.
Consequently
\[
                             \kappa_2\in\Z_2^\times.
                                                               \tag{D.55f}
\]
The use of the global multiplier
$1$ in (D.55e) is essential to
this exact norm equation.

\emph{The primitive boundary rules out every odd fixed character.}
Let $q$ be odd and suppose a nonzero
$\beta$-invariant character
$\chi:F\to C_q$ existed.
Extend it by $\widetilde\chi(s)=0$
to $E_b$, and put
$N=\ker\chi$ and
$E_N=\ker\widetilde\chi=N\rtimes\langle s\rangle$.
Equation (D.55e) gives
$2\widetilde\chi(v)=0$, so $v\in E_N$.
Write $v=ns$ with $n\in N$.

The exact cover and its maps are
those of (D.54g)--(D.54j), now over
$\Z_2$. It has type $(1,q)$.
The deck-average idempotent
$e=q^{-1}\sum_{C_q}d$ commutes
with $\beta$, identifies $eH_1(N,\Z_2)$
with the base torus homology by the
index-$q$ lattice inclusion, and
has boundary augmentation as its
complement.
On the first summand $\beta=-I$
by (D.55d). On the second it is
the identity: the actual $v\in E_N$
fixes each cusp label and centralizes
its cusp generator up to $N$-conjugacy,
exactly as in (D.54g).
Hence
\[
                  I+\beta=0
                       \quad\text{on }H_1(N,\F_2).
                                                               \tag{D.55g}
\]
The same two elements in (D.55e)
then have incompatible homology:
\[
             [v^2]=[n\,\beta(n)]=0,\qquad
             [c^\kappa]=(\kappa_2\bmod2)[c]\ne0.
                                                               \tag{D.55h}
\]
The second class is nonzero because
$c$ is a primitive boundary in
the cover with $q\geq3$ cusps.
This excludes the invariant character.

At every odd prime $q$ the matrix
$A_q(b)$ squares to $I$ and has
determinant $1$.
Its characteristic-zero eigenvalues
are both $1$ or both $-1$.
The first possibility makes it $I$
and supplies a nonzero invariant
character modulo $q$, contrary
to (D.55h).
The second makes it $-I$.
Together with (D.55d), this proves
(D.55a).

\emph{Faithfulness at every prime and the exact abstract list.}
Fix a prime $p$ and a finite
subgroup $G$ of the full torus group.
By D.33 its order divides twenty-four.
A nontrivial kernel of $A_p|_G$
would contain an element of order
two or three.
A positive involution has image
$-I$ by (D.55a).
A negative involution has determinant
$-1$ at every prime by D.43, so
also cannot lie in the kernel.
An order-three element has nonidentity
image at $p=3$ by D.33, and at
every other prime by D.54's
fixed-character obstruction.
The kernel is therefore trivial,
as claimed.

D.54 already removed $C_{12}$
and $D_{24}$ from D.52.
The only additional type in that
allowed list was $\operatorname{Dic}_{12}$.
Its normal subgroup of order three
would be faithfully represented over
$\Q_3$, and an order-four element
inverting it would also have order
four by the just-proved faithfulness.
But the exact field-normalizer
calculation (D.52i)--(D.52k) proves
that every finite-order matrix
inverting this $C_3$ has square $I$:
the alternative would require norm
$-1$ in $\Q_3(\zeta_3)$.
This contradiction excludes
$\operatorname{Dic}_{12}$.
The nine remaining types all have
the explicit realizations of
D.50 and D.52, proving (D.55b).

For an arbitrary order-four $a$,
its square is a positive involution,
so (D.55a) gives (D.55c).
The full multiplier still has square
one by D.50.
At a negative coordinate the two
eigenvalues of $A_q(a)$ are roots
of $X^2+1$ whose product is $-1$,
so they are equal.
The matrix is diagonalizable in
characteristic zero, hence scalar
$i_q I$ or $-i_q I$.
The two-coordinate is positive by
D.50. For an odd $q$ a square
root of $-1$ in $\Q_q$ reduces to
one in $\F_q$, which requires
$q\equiv1\pmod4$.
This proves the stated quartic
restriction without excluding its
remaining scalar case.

The program
\path{code/part_c/torus_cover_norm_d54_d55.py}
also checks the square norm and
primitive-boundary comparison over
$\F_2$ in explicit cyclic odd-degree
covering homology modules.
The full subgroup, covering and
normalizer arguments above give
the all-prime conclusion.
\end{proof}
